% Source pp. 239-253 (Lectures19-25.pdf).
\section{RG and Energy-Momentum Tensor}

Renormalization Group describes the way the scale transformations are realized in
QFT. The scale transformations
\[
x^\mu\to\tilde x^\mu=x^\mu/L
\]
constitute particular class of general coordinate transformations
\[
x^\mu\to\tilde x^\mu=f^\mu(x)\,,
\]
where $f^\mu$ are arbitrary smooth functions. Variations of the action
functional\footnote{The sign minus in this formula corresponds to the relation
$-\mathcal A=iS=i\int dt\,d\xx\,\Lag(x)=\int dx^4\,d\xx\,\Lag(x)=-\int dx^4\,d\xx\,\Lag_E(x)$,
where $\Lag(x)$ is the usual Lagrangian density with $t$ replaced by $-ix_4$, see
Section~7.}
\begin{equation}\label{eq:Aminus}
\mathcal A=-\int d^dx\,\Lag(x)
\end{equation}
under these coordinate transformations in a local field theory is generally
described in terms of its energy-momentum tensor $T^{\mu\nu}(x)$. Therefore this
field, and in particular its trace
\[
\Theta(x):=T^\mu{}_\mu(x)\,,
\]
plays important role in RG theory. In fact, all basic equations of RG theory in
local QFT can be expressed as certain properties of this field.

In the first lecture of this course we have defined the Energy-Momentum tensor (or
EM tensor, for shortness), in a classical field theory, as the Noether's current
associated with the translational symmetry $x^\mu\to x^\mu+a^\mu$,
\[
T^\mu{}_\nu=\frac{\partial\Lag}{\partial(\partial_\mu\varphi)}\,\partial_\nu\varphi-\delta^\mu_\nu\,\Lag\,,
\]
which satisfies the continuity equation
\begin{equation}\label{eq:Tcont}
\partial_\mu T^\mu{}_\nu=0
\end{equation}
as a consequence of the classical equations of motion. Generally, there is an
important ambiguity in the definition of the local EM tensor. One can redefine
$T^\mu{}_\nu$ by adding a total derivative
\[
T^\mu{}_\nu\to T^\mu{}_\nu+\partial_\lambda\psi^{\mu\lambda}{}_\nu\,,\qquad\psi^{\mu\lambda}{}_\nu=-\psi^{\lambda\mu}{}_\nu\,.
\]
The new tensor still satisfies the continuity equation \eqref{eq:Tcont}. By
appropriate choice of $\psi$ the modified EM tensor can always be made symmetric,
$T^{\mu\nu}=T^{\nu\mu}$.

More general way to define the EM tensor is to consider the field theory in the
space with nontrivial metric
\begin{equation}\label{eq:metric}
ds^2=g_{\mu\nu}(x)\,dx^\mu dx^\nu\,.
\end{equation}
In this case the action
\[
\mathcal A[\varphi,g]=-\int d^dx\,\sqrt g\,\Lag(\varphi,\partial_\mu\varphi,g_{\mu\nu})
\]
becomes a functional of the metric tensor $g$ also.

For instance, for generic scalar field theory, in the presence of the metric
background \eqref{eq:metric}, the simplest action is
\[
\mathcal A=\int\sqrt g\,\left[\frac12\,g^{\mu\nu}\,\partial_\mu\varphi\,\partial_\nu\varphi+V(\varphi)\right]\,.
\]
This form is called the ``minimal coupling to gravity''.

In any case, the field $T^{\mu\nu}(x)$ is defined in terms of the variation of the
action with respect to the metric (sign!),
\begin{equation}\label{eq:Tdef}
\delta\mathcal A=-\frac12\int d^dx\,\sqrt g\,T^{\mu\nu}(x)\,\delta g_{\mu\nu}(x)\,.
\end{equation}
Note that $T^{\mu\nu}$ defined this way is automatically symmetric,
\[
T^{\mu\nu}(x)=T^{\nu\mu}(x)\,,
\]
as $g_{\mu\nu}$ is.

Note that for flat space this definition is still ambiguous, as one could couple
the theory to the external metric in many different ways. For example, in the
scalar field theory one could add the ``non-minimal'' terms like
\[
\int d^dx\,\sqrt g\,R\,W(\varphi)\,,\qquad\text{or}\qquad\int d^dx\,\sqrt g\,R_{[\mu\lambda][\nu\rho]}\,W^{[\mu\lambda][\nu\rho]}(\varphi)
\]
where $R$ is the curvature, and $W$ is arbitrary. These terms disappear in the flat
space, however they affect the variation \eqref{eq:Tdef} even when the metric is
set back to the flat Euclidean one. It is possible to show that all such
``non-minimal'' modifications of the action in curved space lead to modifications
of the EM tensor of the form
\begin{equation}\label{eq:Tambig}
T^{\mu\nu}\to T^{\mu\nu}+\partial_\lambda\partial_\rho\,Y^{[\mu\lambda][\nu\rho]}\,,
\end{equation}
where $Y$ is antisymmetric in $[\mu\lambda]$ and in $[\nu\rho]$ but symmetric
w.r.t.\ the interchange of these pairs (this is exactly the symmetry of the
Riemann's curvature tensor $R_{[\mu\lambda][\nu\rho]}$).

This is intrinsic ambiguity which cannot be resolved unless we specify exactly how
the system behaves in curved space. Note that if $Y$ decays sufficiently fast at
$x\to\infty$, such modifications do not affect the notions of energy and momentum,
\[
E=\int d^{d-1}\xx\,T^{00}\,,\qquad P^i=\int d^{d-1}\xx\,T^{0i}\,,
\]
as they add total derivatives to the integrands. If we are interested just in the
theory in flat metric, fixing this ambiguity is a matter of convenience.

Consider now quantum field theory. The correlation functions are defined through
the functional integrals
\begin{equation}\label{eq:corrg}
\bigl\langle\mathcal O_1(x_1)\cdots\mathcal O_n(x_n)\bigr\rangle=Z^{-1}\int[D\varphi]\,\mathcal O_1(x_1)\cdots\mathcal O_n(x_n)\,e^{-\mathcal A[\varphi,g]}\,,
\end{equation}
where for the moment I consider the system in external gravitational field
$g_{\mu\nu}(x)$. Here $\mathcal O(x)$ are some local fields, which can be the field
$\varphi(x)$ itself, or some composite fields built from $\varphi$ and its
derivatives at the point $x$. In the presence of the external metric these
composite fields may depend on this metric in a local way, i.e.\ they may involve
$g_{\mu\nu}(x)$ and its derivatives at the point $x$.

Consider a small variation around the flat metric,
\[
g_{\mu\nu}(x)=\delta_{\mu\nu}+\delta g_{\mu\nu}(x)\,.
\]
For generic variation the correlation functions change in a nontrivial way because
the theory indeed depends on the background metric. However, the change is simple
if we take $\delta g_{\mu\nu}(x)$ corresponding to pure coordinate transformation
\begin{equation}\label{eq:coordtr}
x^\mu\to x^\mu+\varepsilon^\mu(x)
\end{equation}
with infinitesimal $\varepsilon^\mu(x)$; in this case
\[
-\delta g_{\mu\nu}(x)=\partial_\mu\varepsilon_\nu(x)+\partial_\nu\varepsilon_\mu(x)\,.
\]
The coordinate transformation is just a relabeling of the degrees of freedom
$\varphi(x)$, so it amounts to a change of variables in the functional integral
\eqref{eq:corrg}. Therefore the numerical value of this integral does not change.
This leads to the identity
\begin{multline}\label{eq:Ward}
0=\sum_{i=1}^n\bigl\langle\mathcal O_1(x_1)\cdots\delta_\varepsilon\mathcal O_i(x_i)\cdots\mathcal O_n(x_n)\bigr\rangle\\
+\int d^dx\,\partial_\mu\varepsilon_\nu(x)\,\bigl\langle T^{\mu\nu}(x)\,\mathcal O_1(x_1)\cdots\mathcal O_n(x_n)\bigr\rangle\,,
\end{multline}
where $\delta_\varepsilon\mathcal O_i(x)$ stand for local variations of the fields
$\mathcal O_i(x_i)$ under the coordinate transformation \eqref{eq:coordtr}. The
identities of this type are generally known as the \emph{Ward identities}.

What can be said about the variations $\delta_\varepsilon\mathcal O_i(x)$ on
general grounds? These again are local fields which depend on $\varepsilon^\mu(x)$
\begin{itemize}
\item[(i)] linearly, i.e.\
$\delta_{\varepsilon_1+\varepsilon_2}\mathcal O=\delta_{\varepsilon_1}\mathcal O+\delta_{\varepsilon_2}\mathcal O$,
and
\item[(ii)] locally, i.e.\ $\delta_\varepsilon\mathcal O(x)$ depends on
$\varepsilon^\mu(x)$ and its derivatives taken at the point $x$.
\end{itemize}
In fact, the part of $\delta_\varepsilon\mathcal O(x)$ which contains just
$\varepsilon^\mu(x)$, with no derivatives, is completely fixed,
\[
\delta_\varepsilon\mathcal O(x)=\varepsilon^\mu(x)\,\partial_\mu\mathcal O(x)+\text{terms with }\partial_\nu\varepsilon^\mu(x)\text{ and higher derivatives}\,,
\]
because for constant $\varepsilon^\mu(x)=\varepsilon^\mu$ we have
\[
\mathcal O(x)\to\mathcal O(x+\varepsilon)=\mathcal O(x)+\varepsilon^\mu\,\partial_\mu\mathcal O(x)\,.
\]
Moreover, the part which contains the first derivatives $\partial_\nu\varepsilon^\mu(x)$
is also severely restricted. For instance, if $\varphi(x)$ is a scalar field, and we
assume that it transforms independently of other fields, we can write
\[
\delta_\varepsilon\varphi(x)=\varepsilon^\mu(x)\,\partial_\mu\varphi(x)+\frac Dd\,(\partial_\mu\varepsilon^\mu)\,\varphi(x)\,,
\]
where $d$ is the dimensionality of the space, and the parameter $D$ is related to
the dimension of $\varphi$. This generalizes to more complicated scalar composite
fields $\mathcal O_\alpha$ as follows
\[
\delta_\varepsilon\mathcal O_\alpha(x)=\varepsilon^\mu(x)\,\partial_\mu\mathcal O_\alpha(x)+\frac{D_\alpha{}^\beta}d\,(\partial_\mu\varepsilon^\mu)\,\mathcal O_\beta(x)+\text{terms with }\partial_\nu\partial_\lambda\varepsilon^\mu\,,\ \text{etc.}
\]
For tensor fields the terms involving local rotations
\[
\omega_{\mu\nu}(x)=\partial_\mu\varepsilon_\nu(x)-\partial_\nu\varepsilon_\mu(x)
\]
will also appear.

Going back to the Ward Identity \eqref{eq:Ward}, we observe that integral there can
be transformed by parts, hence we have
\begin{multline}\label{eq:Ward2}
0=\sum_{i=1}^n\bigl\langle\mathcal O_1(x_1)\cdots\delta_\varepsilon\mathcal O_i(x_i)\cdots\mathcal O_n(x_n)\bigr\rangle\\
-\int d^dx\,\varepsilon_\nu(x)\,\bigl\langle\partial_\mu T^{\mu\nu}(x)\,\mathcal O_1(x_1)\cdots\mathcal O_n(x_n)\bigr\rangle\,.
\end{multline}
As the first term here depends only on $\varepsilon(x)$ (and its derivatives) at the
points $x_1,\dots,x_n$ only, we conclude that
\[
\bigl\langle\partial_\mu T^{\mu\nu}(x)\,\mathcal O_1(x_1)\cdots\mathcal O_n(x_n)\bigr\rangle=0\qquad\text{for}\quad x\neq x_1,\dots,x_n\,,
\]
i.e.\ in quantum field theory the continuity equation \eqref{eq:Tcont} is satisfied
in the weak sense,
\[
\partial_\mu T^{\mu\nu}(x)\simeq0\,,
\]
exactly as the equations of motion do.

Now let us consider particular case of the coordinate transformations, the scale
transformations (also called the dilations)
\[
x^\mu\to x^\mu+\varepsilon\,x^\mu\,.
\]
In this case
\[
\partial_\mu\varepsilon_\nu(x)=\delta_{\mu\nu}\,\varepsilon\,,
\]
while all higher derivatives of $\varepsilon(x)$ vanish, and so for scalar fields
$\mathcal O_\alpha$ we have
\[
\delta_\varepsilon\mathcal O_\alpha(x)=\varepsilon\,(\mathcal D\mathcal O)_\alpha(x)\,,
\]
where by definition
\begin{equation}\label{eq:DO}
(\mathcal D\mathcal O)_\alpha(x)=x^\mu\frac\partial{\partial x^\mu}\,\mathcal O_\alpha(x)+D_\alpha{}^\beta\,\mathcal O_\beta(x)\,.
\end{equation}
Then the Ward identity \eqref{eq:Ward} reduces to
\begin{equation}\label{eq:WardTheta}
\sum_{i=1}^n\bigl\langle\mathcal O_{\alpha_1}(x_1)\cdots(\mathcal D\mathcal O)_{\alpha_i}(x_i)\cdots\mathcal O_{\alpha_n}(x_n)\bigr\rangle
=-\int d^dx\,\bigl\langle\Theta(x)\,\mathcal O_{\alpha_1}(x_1)\cdots\mathcal O_{\alpha_n}(x_n)\bigr\rangle
\end{equation}
where
\[
\Theta(x)=T^\mu{}_\mu(x)
\]
is the trace component of the EM tensor.

Now let us assume that our QFT depends on some number $N$ of parameters (the
``coupling constants'') $\lambda^k$, $k=1,\dots,N$. These parameters enter the
action $\mathcal A$. The derivatives
\[
\Phi_k(x)=-\frac\partial{\partial\lambda^k}\,\Lag(x)
\]
are in general some scalar composite fields. Then\footnote{Of course, in general
the fields $\mathcal O_\alpha(x)$ may depend on the couplings $\lambda$. If such
dependence is present, the derivatives of $\mathcal O_\alpha$ would add to the
matrix $D_\alpha{}^\beta$, but otherwise would not change our argument. Here I ignore
this contribution for simplicity.}
\[
\frac\partial{\partial\lambda^k}\bigl\langle\mathcal O_{\alpha_1}(x_1)\cdots\mathcal O_{\alpha_n}(x_n)\bigr\rangle
=-\int d^dx\,\bigl\langle\Phi_k(x)\,\mathcal O_{\alpha_1}(x_1)\cdots\mathcal O_{\alpha_n}(x_n)\bigr\rangle
\]
If we assume that $\Theta(x)$ can be expanded in terms of $\Phi_k(x)$,
\begin{equation}\label{eq:ThetaBeta}
\Theta(x)=-\sum_{k=1}^N\beta^k(\lambda)\,\Phi_k(x)\,,
\end{equation}
one arrives again at the Callan-Symanzik equation, i.e.
\begin{multline}\label{eq:CSTheta}
\sum_{i=1}^n\bigl\langle\mathcal O_{\alpha_1}(x_1)\cdots(\mathcal D\mathcal O)_{\alpha_i}(x_i)\cdots\mathcal O_{\alpha_n}(x_n)\bigr\rangle\\
+\sum_{k=1}^N\beta^k(\lambda)\,\frac\partial{\partial\lambda^k}\bigl\langle\mathcal O_{\alpha_1}(x_1)\cdots\mathcal O_{\alpha_n}(x_n)\bigr\rangle=0\,,
\end{multline}
where the coefficients $D_\alpha{}^\beta(\lambda)$ in the operator $\mathcal D$ in
\eqref{eq:DO} can of course depend on the parameters $\lambda$. The possibility to
expand the field $\Theta(x)$ in terms of finitely many ``basic'' fields
$\Phi_k(x)$ is the statement of renormalizability of the field theory.

The energy-momentum tensor $T^{\mu\nu}(x)$ is special composite field. As we know,
in interacting field theory composite fields in general may require infinite
renormalizations. However, the above derivation of the Ward identity
\eqref{eq:Ward} suggests that in the case of the EM tensor no infinite
renormalization is needed. Indeed, the correlation functions in \eqref{eq:Ward}
are understood as the correlation functions of fully renormalized fields
$\mathcal O_\alpha$, they remain finite when the regularization is removed.
Therefore the EM tensor defined through the variation \eqref{eq:Tdef} must have
finite correlation functions without additional renormalization. It is possible to
check this statement in renormalized perturbation theory in $d=4$ $\varphi^4$
theory.

\subsection*{Conformal invariance at the fixed-point}

If $\lambda^k=\lambda^k_*$ such that
\[
\beta^k(\lambda_*)=0\,,
\]
the trace of the EM tensor vanishes in view of \eqref{eq:ThetaBeta}
\begin{equation}\label{eq:Theta0}
\Theta(x)=0\,.
\end{equation}

\begin{remark}
The above statement ignores important subtlety. The derivation of the
Callan-Symanzik equation \eqref{eq:CSTheta} relies in fact on the integrated version
of \eqref{eq:ThetaBeta}. Locally, Eq.~\eqref{eq:ThetaBeta} may hold only up to total
derivatives, i.e.\ the terms of the form $\partial_\mu K^\mu(x)$, with $K^\mu(x)$
being some local field. Thus, vanishing of the beta-functions only implies that
$\Theta(x)=\partial_\mu K^\mu(x)$. Unless $K^\mu$ is itself a total derivative
(i.e.\ $\Theta=\partial_\mu\partial_\nu L$, in which case the ambiguity
\eqref{eq:Tambig} can be used to redefine $T_{\mu\nu}$ to eliminate its trace),
presence of such a term would spoil the argument for conformal invariance. In some
cases (in particular, in $d=2$ unitary quantum field theories) this possibility can
be excluded. Generally it is not excluded, and therefore vanishing of the
beta-functions does not guarantee conformal invariance. Moreover, in general,
counterexamples can be found. For more on this subtlety, see \emph{J.~Polchinski,
Renormalization and effective Lagrangians, Nucl.\ Phys.\ B231, 269 (1984)}. In the
following discussion I assume that the fixed point in question admits the EM tensor
which satisfies \eqref{eq:Theta0}.
\end{remark}

The correlation functions of this fixed-point theory enjoy scale invariance with
the dimensions of the fields $\mathcal O_\alpha$ determined by the eigenvalues of
the matrix
\[
D_\alpha{}^\beta=D_\alpha{}^\beta(\lambda_*)\,.
\]
In fact, vanishing of the trace component $\Theta(x)$ in the fixed-point theory
guarantees much bigger symmetry. Namely, let $v^\mu(x)$ be any vector field which
satisfies the equation
\begin{equation}\label{eq:CKV}
\partial_\mu v_\nu(x)+\partial_\nu v_\mu(x)=\rho(x)\,\delta_{\mu\nu}=\frac2d\,\delta_{\mu\nu}\,\bigl(\partial_\lambda v^\lambda(x)\bigr)\,.
\end{equation}
Consider the coordinate transformation
\begin{equation}\label{eq:confeps}
x^\mu\to x^\mu+\varepsilon^\mu(x)\qquad\text{with}\qquad\varepsilon^\mu(x)=\varepsilon\,v^\mu(x)\,,
\end{equation}
where $\varepsilon$ is infinitesimal parameter. For such $\varepsilon^\mu$
\[
\partial_\mu\varepsilon_\nu\,T^{\mu\nu}=\frac\varepsilon2\,\rho(x)\,\Theta(x)=0\,,
\]
and the second term in the Ward identity \eqref{eq:Ward} disappears, i.e.
\begin{equation}\label{eq:WardConf}
\sum_{i=1}^n\bigl\langle\mathcal O_1(x_1)\cdots\delta_\varepsilon\mathcal O_i(x_i)\cdots\mathcal O_n(x_n)\bigr\rangle=0\,.
\end{equation}
We obtain the closed equation
\begin{equation}\label{eq:WardConf2}
\sum_{i=1}^n\bigl\langle\Phi_{\alpha_1}(x_1)\cdots\mathcal D_v\Phi_{\alpha_i}(x_i)\cdots\Phi_{\alpha_n}(x_n)\bigr\rangle=0\,,
\end{equation}
where
\[
\mathcal D_v\Phi_\alpha(x)=\left(v^\mu(x)\,\frac\partial{\partial x^\mu}+\frac{D_\alpha}d\,\bigl(\partial_\mu v^\mu(x)\bigr)\right)\,\Phi_\alpha(x)\,,
\]
and $\Phi_\alpha(x)$ are the fields which diagonalize the matrix $D_\alpha{}^\beta$,
and $D_\alpha$ are corresponding eigenvalues.

The coordinate transformations \eqref{eq:confeps} which satisfy the equation
\eqref{eq:CKV} are called the \emph{conformal transformations}. It is possible to
find all solutions of this equation. The situation turns out to be rather different
in the cases $d>2$ and $d=2$.

If $d>2$ the general solution has the form\footnote{The Lie algebra of these vector
fields is isomorphic to $so(d+1,1)$, the algebra of infinitesimal Lorentz
transformations of $d+2$-dimensional space-time.}
\[
v^\mu(x)=a^\mu+\omega^{\mu\nu}x^\nu+l\,x^\mu+\bigl(b^\mu x^2-2x^\mu(bx)\bigr)\,,
\]
where $\omega^{\mu\nu}=-\omega^{\nu\mu}$. The parameters $a,\omega,l,b$ are
interpreted as follows:
\begin{center}
\begin{tabular}{lll}
$a^\mu$ & -- & translations\\
$\omega^{\mu\nu}$ & -- & rotations\\
$l$ & -- & dilations\\
$b^\mu$ & -- & special conformal transformations
\end{tabular}
\end{center}
As the first three are very familiar, I will present only the finite version of the
special conformal transformations,
\[
x^\mu\to\tilde x^\mu=\frac{x^\mu+B^\mu x^2}{1+2(Bx)+B^2x^2}\,,
\]
where $B^\mu$ is an arbitrary vector.

Invariance of the fixed-point theory w.r.t.\ the special conformal transformations
provides further limitations on the correlation functions. Consider the Ward
identity \eqref{eq:WardConf2} with
\[
v^\mu(x)=b^\mu x^2-2x^\mu(bx)\,.
\]
In this case
\[
\mathcal D_v\Phi_\alpha(x)=\left[\bigl(b^\mu x^2-2x^\mu(bx)\bigr)\frac\partial{\partial x^\mu}-2D_\alpha\,(bx)\right]\,\Phi_\alpha(x)\,,
\]
and \eqref{eq:WardConf2} becomes nontrivial closed differential equation for the
correlation functions. In fact, this equation completely determines one-, two- and
three-point correlation functions, up to overall constant factors (this is left as
an exercise).

If $d=2$ the equation
\begin{equation}\label{eq:CKV2}
\partial_\mu v_\nu+\partial_\nu v_\mu-\delta_{\mu\nu}\,(\partial_\lambda v^\lambda)=0
\end{equation}
has infinitely many solutions. This is very easy to see by going to the complex
coordinates in $d=2$ space,
\[
z=x^1+ix^2\,,\qquad\bar z=x^1-ix^2\,.
\]
Then \eqref{eq:CKV2} leads to two equations
\[
\partial_{\bar z}v^z=0\,,\qquad\partial_zv^{\bar z}=0\,,
\]
which shows that $v^z(z)$ ($v^{\bar z}(\bar z)$) can be arbitrary holomorphic
(anti-holomorphic) function. This rich symmetry has important consequences for
the structure of $d=2$ conformal field theories.

Scale invariance of the fixed-point theory allows one to establish some important
properties of QFT associated with its unstable manifold.

\subsection*{Scale invariance and density of states}

Important conclusion can be made about its density of states. The density of
states is directly related to the equilibrium state entropy $S(E)$ of the system.
The form of $S(E)$ in a scale invariant quantum field theory can be established
directly by general arguments. However, I will use this opportunity to display some
general features of QFT at finite temperature.

The functional integral of $d$-dimensional Euclidean quantum field theory
\[
\langle\dots\rangle=Z^{-1}\int D[\varphi]\,(\dots)\,e^{-\mathcal A[\varphi]}
\]
has formal interpretation as expectation value over a thermal equilibrium ensemble
in \emph{classical} statistical mechanics in $d$ spatial dimensions. We have used
this interpretation many times already. Here I would like to discuss Euclidean
functional integral representation of \emph{quantum} statistical mechanics of
fields in $d$ dimensional space-time (i.e.\ $d-1$ dimensional space).

Recall that the density matrix of thermal equilibrium state in quantum statistics
has the general form\footnote{The normalization factor $Z(\beta)$ here is different
from the factor $Z$ in e.g.\ Eq.~\eqref{eq:corrdef}. $Z(\beta)$ goes to $Z$ in the
low temperature limit $\beta\to\infty$ (see below).}
\begin{equation}\label{eq:density}
\hat P=Z^{-1}(\beta)\,e^{-\beta\hat H}\,,
\end{equation}
where $\beta=(kT)^{-1}$. The exponential here coincides with the time evolution
operator at pure imaginary value of time $t=-i\hbar\beta$ (I set $\hbar=1$
henceforth). This form is still valid if one is dealing with the equilibrium state
of quantum fields. The imaginary-time evolution operator has representation as the
functional integral,
\begin{equation}\label{eq:PIbeta}
\bigl\langle\varphi_f(\xx)\big|e^{-\beta\hat H}\big|\varphi_i(\xx)\bigr\rangle=\int D[\varphi(\xx,\tau)]\Big|_{\substack{\varphi(\xx,0)=\varphi_i(\xx)\\\varphi(\xx,\beta)=\varphi_f(\xx)}}\,e^{-\mathcal A[\varphi]}
\end{equation}
Here $\xx$ are $d-1$ dimensional ``spatial'' coordinates, and $\tau$ is the
imaginary time (so that $x=(\xx,\tau)$). The integration is over the functions
$\varphi(x)=\varphi(\xx,\tau)$ in the slab $0\le\tau\le\beta$, with the boundary
conditions
\[
\varphi(\xx,0)=\varphi_i(\xx)\,,\qquad\varphi(\xx,\beta)=\varphi_f(\xx)\,.
\]
The partition sum of the quantum system at the finite temperature is the trace
\begin{equation}\label{eq:Zbeta}
Z(\beta)=\tr\bigl(e^{-\beta\hat H}\bigr)\,,
\end{equation}
i.e.\ the sum of its diagonal matrix elements
\begin{equation}\label{eq:Ztrace}
\int D[\varphi(\xx)]\,\bigl\langle\varphi(\xx)\big|e^{-\beta\hat H}\big|\varphi(\xx)\bigr\rangle\,.
\end{equation}
If the representation \eqref{eq:PIbeta} is used, the integration over
$\varphi(\xx)$ in \eqref{eq:Ztrace} combines with the integral $D[\varphi(x)]$ in
\eqref{eq:PIbeta} to produce integral over all fields $\varphi(\xx,\tau)$, with
$0\le\tau\le\beta$, satisfying the periodic boundary condition
\begin{equation}\label{eq:periodic}
\varphi(\xx,0)=\varphi(\xx,\beta)\,.
\end{equation}
Thus, we have
\begin{equation}\label{eq:Zcyl}
Z(\beta)=\int D[\varphi(\xx,\tau)]\big|_{\varphi(\xx,0)=\varphi(\xx,\beta)}\,e^{-\mathcal A[\varphi]}\,.
\end{equation}
Geometrically, this is the functional integral over the functions $\varphi(x)$ on a
$d$-dimensional cylinder $\R^{d-1}\times S^1$, with $\tau$ being the coordinate on
the circle of circumference $\beta$, $\tau\sim\tau+\beta$. Of course this is nothing
but the Matsubara representation specialized to the quantum field theory at hand.

Now, let us change interpretation of the functional integral \eqref{eq:Zcyl}. We
can take one of the coordinates on $\R^{d-1}$, say $x_1$, and claim it to be the
imaginary time. The coordinate $\tau$ will be treated as one of the spatial
coordinates. In this interpretation the imaginary time spans infinite range,
$-\infty<x_1<\infty$, while the ``space'' $\R^{d-2}\times S^1$ is the $d-1$
dimensional cylinder, with the coordinates $(\xx_\parallel,x_\perp)$. Here
$\xx_\parallel=(x_2,\dots,x_{d-1})$, and $x_\perp=\tau$. The latter is
compactified,
\[
x_\perp\sim x_\perp+\beta\,.
\]
Since now the imaginary time spans large (in the limit, infinite) range,
\eqref{eq:Zcyl} is dominated by the contribution from the ground state,
\begin{equation}\label{eq:ZE0}
Z(\beta)\sim e^{-L\,E_0(\beta)}\,,
\end{equation}
where $L\to\infty$ is the range of $x_1$, and $E_0(\beta)$ is the ground state
energy of the quantum field in the spatial geometry $\R^{d-2}\times S^1$.

This geometric setting is similar to that we discussed in the context of the
Casimir effect, with $\beta$ replacing the separation between the plates $a$. The
difference is that in the Casimir case we had a fixed boundary condition
$\varphi=0$ at the plates, which now is replaced by the periodic boundary condition
\eqref{eq:periodic}. For the free massless field it is easy to repeat the
calculation of the Casimir energy (I recommend it as an Exercise), with the result
\begin{equation}\label{eq:E0beta}
E_0(\beta)=-\frac{C\,V_\parallel}{\beta^{d-1}}\,,
\end{equation}
where $C$ is some dimensionless ($d$-dependent) constant, and
\[
V_\parallel=\int d^{d-2}\xx_\parallel
\]
is the volume of the $\R^{d-2}$ component of the space. In fact, apart from the
value of the constant, the form \eqref{eq:E0beta} is dictated by scale invariance
of the theory. Indeed, the ground state energy is an extensive quantity, which at
large $V_\parallel$ must be proportional to this volume. Since the theory has no
intrinsic scale, the $\beta$-dependence then follows from the dimensional analysis.
This simple argument shows that the form \eqref{eq:E0beta} applies to any
scale-invariant field theory, id est fixed point field theory. The constant $C$ of
course would be different for different fixed points.

Returning to original interpretation (the quantum field at finite temperature
$kT=\beta^{-1}$), and using \eqref{eq:ZE0} we have for the free energy of a
scale-invariant field theory
\begin{equation}\label{eq:Fbeta}
\beta F(\beta)=-\log\bigl(Z(\beta)\bigr)=-\frac{C\,V_{d-1}}{\beta^{d-1}}\,,
\end{equation}
where $V_{d-1}=V_\parallel\times L$ is the volume of the space. Of course, one could
have established the form \eqref{eq:Fbeta} directly, by dimensional analysis (the
free energy $F$, being extensive characteristic, must be $\sim V_{d-1}$, whereas
the quantity \eqref{eq:Fbeta} is dimensionless), but I believe our detour through
the Casimir energy was instructive.

The finite temperature partition function is closely related to the energy density
of states. We have
\[
Z(\beta)=\int_0^\infty e^{-\beta E}\,\rho(E)\,dE\,,
\]
where $\rho(E)\,\Delta E$ is the number of states in the energy interval between
$E$ and $E+\Delta E$. Mathematically, $Z(\beta)$ is the Laplace transform of
$\rho(E)$. The inverse transform is
\[
\rho(E)=\int_{-i\infty}^{+i\infty}e^{\beta E}\,Z(\beta)\,\frac{d\beta}{2\pi i}
=\int_{-i\infty}^{+i\infty}e^{\beta E}\,e^{\frac{C\,V_{d-1}}{\beta^{d-1}}}\,\frac{d\beta}{2\pi i}\,.
\]
At large $V_{d-1}$ the integral is dominated by the saddle point at
\[
\beta=\beta_0(E):=\left(\frac{C\,(d-1)\,V_{d-1}}E\right)^{1/d}\,,
\]
and we find
\begin{equation}\label{eq:rhoE}
\rho(E)\sim e^{S(E)}\,,
\end{equation}
where
\begin{equation}\label{eq:SE}
S(E)=\tilde C\,V_{d-1}^{1/d}\,E^{1-1/d}
\end{equation}
where $\tilde C$ is another constant, related to $C$ in
\eqref{eq:E0beta}.\footnote{$\tilde C=\frac d{d-1}\,\bigl(C\,(d-1)\bigr)^{1/d}$.}

By the definition, the quantity $S(E)$ in \eqref{eq:rhoE} coincides with the
equilibrium state entropy of the system at the temperature $kT=1/\beta_0(E)$. The
equation \eqref{eq:SE} expresses the universal scale invariant form of the entropy.
Indeed, the entropy is an extensive property which must have the form
$V_{d-1}\,\sigma$, where the entropy density $\sigma=\sigma(\varepsilon)$ can depend
on the energy density $\varepsilon:=E/V_{d-1}$. On the other hand $S(E)$, being a
dimensionless quantity, can only depend on the dimensionless combination
$V_{d-1}E^{d-1}$. Eq.~\eqref{eq:SE} is the only form that satisfies these
conditions. Thus, the asymptotic density of states \eqref{eq:rhoE}, \eqref{eq:SE}
can be deduced directly, from the scale invariance of the fixed point theory.
Nonetheless, our previous discussion is instructive as it displays some basics of
the finite-temperature formalism of QFT.

Note that in any dimensionality $S(E)$ grows slower than the exponential $\sim E$.
This observation is useful in comparing to the string theories, whose typical
energy density is the exponential.\footnote{Stringy density of states
$S(E)\simeq E/kT_H$ has manifestation in the so called \emph{Hagedorn transition}
at $T=T_H$.}

It is straightforward to argue that the energy density \eqref{eq:rhoE},
\eqref{eq:SE} is valid asymptotically, at large $E$, in any quantum field theory
originating at certain critical fixed point (i.e.\ associated with
$U(\mathcal A_*)$). And if we adopt the Wilsonian idea of quantum field theories as
RG flows on the unstable manifolds of fixed points, the above equation \eqref{eq:SE}
becomes general limitation on the asymptotic density of states in quantum field
theory.

Let me make here a remark on the status of this idea. In our discussion of the
Renormalization Group we (following Wilson) found it natural to identify local field
theories with the special class of RG flows in $\Sigma$, the flows that can be
``integrated backward'' without limit. A critical fixed point $\mathcal A_*$ gives
simple mechanism for generating such flows. There can be several fixed points, or
even continuous manifold of fixed points.\footnote{About limit cycles,
$C$-theorems, etc.} In any case, the above conclusion about the high energy density
of states, which follows from the scale invariance of the fixed point theory,
remains valid.

However, another, somewhat more intricate behavior is conceivable. Let me briefly
discuss it, just in the way of entertainment. To start with, let us imagine a slice
of the critical surface passing through a fixed point $\mathcal A_*$. All RG
trajectories in some domain around $\mathcal A_*$ flow to $\mathcal A_*$ at
$l\to+\infty$, see figure~\ref{fig:nestedfp}. Let's also assume that there is
another fixed point $\mathcal A'_*$ lying in the same critical surface, such that
$\mathcal A_*\in U(\mathcal A'_*)$. Furthermore, assume there is yet another fixed
point $\mathcal A''_*\in U(\mathcal A'_*)$.
\begin{figure}[ht]
\centering\figfile{nestedfp}
\caption{An RG trajectory starting at $\mathcal A''_*$, passing close to
$\mathcal A'_*$, and ending at $\mathcal A_*$.}
\label{fig:nestedfp}
\end{figure}

This figure shows an RG trajectory which originates at $\mathcal A''_*$, then passes
very close to $\mathcal A'_*$, finally converging to $\mathcal A_*$. Starting from
any point on this trajectory, one can ``integrate backward'' without limit, i.e.\
it defines local QFT. While the large scale properties of this QFT are controlled by
the ``infrared'' fixed point $\mathcal A_*$, its behavior at shorter scales is more
intricate. When traced backwards, the RG trajectory stays long ``time'' near the
fixed point $\mathcal A'_*$, before departing towards $\mathcal A''_*$ at
$l\to-\infty$. The associated QFT has its asymptotic short distance behavior
determined by the ``ultraviolet'' fixed point $\mathcal A''_*$, but in some wide
range of intermediate scales it has scaling properties associated with
$\mathcal A'_*$.

One can imagine yet more complicated sequence of nested fixed points, giving rise to
QFT with many different intermediate regimes corresponding to the trajectory passing
close to consecutive fixed points. Ultimately, then, a flow is conceivable which
passes successively near infinitely many fixed points. While such trajectory can be
integrated backward with no limit, its short distance behavior is not controlled by
any given fixed point.
